\documentclass[11pt]{article}
\usepackage[a4paper,margin=27mm]{geometry}
\usepackage[T1]{fontenc}
\usepackage{lmodern,microtype,amsmath,amssymb,amsthm,mathtools,enumitem}
\usepackage[hidelinks]{hyperref}
\newtheorem{theorem}{Theorem}
\newtheorem{lemma}[theorem]{Lemma}
\newtheorem{corollary}[theorem]{Corollary}
\theoremstyle{remark}\newtheorem{remark}[theorem]{Remark}
\newcommand{\Dscr}{\mathcal D}
\newcommand{\ess}{\operatorname{ess}}
\title{Dense Thresholds and Debt Reduction}
\author{}\date{10 October 2026}
\begin{document}\maketitle

\noindent
This note strengthens \emph{Independence Chains} (IC). It does not settle
the general conjecture. It proves the following restrictions on a proposed
counterexample with two ordinary additive worlds:
\begin{enumerate}[leftmargin=*,itemsep=2pt]
\item The thresholds of IC's useful full-support selections are dense in
\emph{all} of $(0,\infty)$, not merely arbitrarily close to zero.
\item Both threshold functions $F,G$ must have jumps in every open interval.
In particular, a single interval of continuity, anywhere, suffices for
impossibility. Even a single constant interval suffices without any further
condition on the penalties.
\item A positive lower bound on $c_3/c_2$ in a strip $t_{12}<h$, or a finite
upper bound in a strip $t_{13}<h$, suffices for impossibility. In particular, a global upper
bound alone, or a global positive lower bound alone, suffices. No continuity
assumption on $F$ or $G$ is needed.
\item For a separable ratio $\rho=f(t_{12})/h(t_{13})$ with positive
nondecreasing factors, continuity of either factor suffices. More generally,
it suffices that the factor's multiplicative jump ratios are bounded
for all sufficiently large arguments.
\end{enumerate}
The earlier version was independently reviewed in \emph{Closed Choice
Forces Fineness} (CCF). That review found an invalid placement of the
silent reservoir in the original debt-contraction argument. The local
cycle below states its geometric hypotheses explicitly, and the following
flexible-credit lemma writes out CCF's proposed repair. The expanded repair
has been checked by the author, not independently refereed. The earlier
notes are used only for the facts explicitly listed below.

\section{Setting and the facts used from IC}

There are three worlds. Smaller loss is better. Models are compared modulo
null sets. Guidance ranks good advice, silence, bad advice and contradictory
advice in that order. We assume strict guidance monotonicity (G), independence
under replacement of common advice (IB), and closed competitive selection (S).
Thus $C_p$ is coherent and undominated, competitive on $\operatorname{supp}p$,
and $\{p:C_p\preceq_iX\}$ is closed for every model $X$.
At worlds $2,3$, the rankings are represented by ordinary additive losses
$L_2,L_3$ with strictly ordered local penalties and an integrable envelope.
World $1$ can have any preorder satisfying (G) and (IB).

Assume, for a contradiction, that the vacuous model $V$ is undominated.
Represent binary problems by $H=\{g_1>0\}$, with its finite Gaussian measure
$\mu$, and put
\[
 P_1=\{g_2<0,g_3<0\},\quad P_2=\{g_2<0<g_3\},\quad
 P_3=\{g_2>0>g_3\}.
\]
The option $o_K$ is good exactly at worlds in $K$. At world $1$, accepting
$g$ gives \emph{credit}, accepting $-g$ gives \emph{debt}, and neither gives
silence. Thus the credit/debt options on $P_1,P_2,P_3$ are respectively
$(o_1,o_{23})$, $(o_{13},o_2)$ and $(o_{12},o_3)$.
Write
\[
 x=t_{12}=g_1/|g_2|,\qquad y=t_{13}=g_1/|g_3|,\qquad
 t=t_{23}=|g_2|/|g_3|
\]
where applicable. All relevant ratio intervals and rectangles have positive
measure. At worlds $k=2,3$, normalize good advice to cost zero, and write
$c_k,b_k$ for the silence and bad costs. Put $u=b_3-c_3$ and
$\rho=c_3/c_2$ on $P_1$. All these densities are positive and integrable.
For a density $v$, write $v(A)=\int_Av\,d\mu$.

Here are the facts from IC and its stated predecessors that we use.
\begin{enumerate}[label=(\roman*),leftmargin=*,itemsep=3pt]
\item Every model $X$ has a selected coherent $C$ with $C\preceq_iX$ at
all worlds (the ordinal KKM lemma). An undominated $V$ is competitive.
\item There are finite positive nondecreasing functions $F,G$ such that
\[
 c_2=F(t)u,\qquad b_2-c_2=G(t)c_3\quad\text{on }P_2\cup P_3,
 \qquad \sup F\le\inf G.
\]
If $\{F=\kappa\}$ and $\{G=\kappa\}$ both have positive Lebesgue measure,
$V$ is dominated (IC's tied-level lemma).
\item All the edge models $E_2(a)$ and $E_3(b)$ occur as selections.
$E_2(a)$ gives debt on $(P_1\cup P_2)\cap\{x<a\}$ and silence on the
remainder of $P_1\cup P_2$, and gives credit on $P_3$.
$E_3(b)$ gives debt on $(P_1\cup P_3)\cap\{y<b\}$ and silence on the
remainder of $P_1\cup P_3$, and gives credit on $P_2$.
For finite positive parameters they are competitive on their respective
two-world supports.
\item A full-support selection cannot give both $o_{13}$ and $o_2$ on
positive parts of $P_2$ (IC's no-both lemma). A selection on $(e_2,e_3)$
cannot give both $o_2$ on $P_2$ and $o_3$ on $P_3$ (IC's weights lemma).
\item A \emph{good base} is a full-support selection $C=\Dscr_K$ such that
$e_2\in K$, $\min_Km_2>0$, and $K$ is contained in neither
$\{m_1=0\}$ nor $\{m_3=0\}$. Its parameters are
\[
 \tau=\max_K\frac{m_3}{m_2}>0,\qquad
 \tau_2=\min_{m\in K:m_1>0}\frac{m_2}{m_1}>0.
\]
It gives no credit on $P_1\cup P_2$ and no debt on $P_3$.
Its debt on $P_2$ contains
\[
 \{t>\tau,\ x<\tau_2(1-\tau/t)\}.
\]
Its silence on $P_3$ lies in $\{t\le\tau\}$ and meets every band
$\{\tau-\delta<t<\tau\}$ in a nonempty open set.
Its silence on $P_1$ has positive measure and lies in $\{x\ge\tau_2\}$.
\item Closedness on $(e_1,e_3)$ forbids a fixed positive credit set
$S\subset P_3\cap\{y\ge b\}$ from being financed by the entire debt band
\[
 B_n=(P_1\cup P_3)\cap\{b-1/n\le y<b\}
\]
for every sufficiently large $n$: one cannot have
$V\preceq_1V+o_{12}(S)+\mathrm{debt}(B_n)$ for all such $n$
(IC's closing lemma).
\end{enumerate}
To use competitiveness in a chain, make a change which strictly decreases
both $L_2,L_3$ at a full-support base, or decreases $L_k$ at a base competitive
on $\{1,k\}$. The base must rank itself no worse than the changed model at
world $1$. By (IB), this comparison transfers to every context with the same
advice on the change set. Finite chains compose by transitivity.

\section{Good bases near every threshold}

For $\theta>0$, let $B_\theta$ have credal set
\[
 K_\theta=\operatorname{conv}\left\{e_2,
       \left(0,\frac1{1+\theta},\frac\theta{1+\theta}\right)\right\}.
\]
It gives debt on $P_1$, debt on $P_2\cap\{t>\theta\}$, credit on
$P_3\cap\{t>\theta\}$, and silence on the remaining mixed problems.
The symbol $B_\theta$ specifies a model; we do not assume that it is selected.

\begin{lemma}\label{lem:bayes}
For every finite $\theta>0$, some finite $\kappa>0$ makes $B_\theta$
a minimizer of $L_2+\kappa L_3$ over all models. Every coherent minimizer
is a model $B_s$, with $s\in[0,\infty]$. Among these minimizers,
the two inequalities $L_k(C)\le L_k(B_\theta)$, $k=2,3$, force
$C=B_\theta$ modulo null sets.
\end{lemma}
\begin{proof}
Take $\kappa=F(\theta+)$. Then $F(t)\le\kappa$ below $\theta$ and
$F(t)\ge\kappa$ above it. Also $G>\kappa$ almost everywhere. Indeed this is
immediate if $\kappa<\inf G$. Otherwise $\kappa=\sup F=\inf G$, and $F$
equals $\kappa$ on $(\theta,\infty)$; the tied-level lemma forces
$\{G=\kappa\}$ to be null.

On $P_1$, $o_{23}$ uniquely minimizes the weighted loss. On $P_2$ and
$P_3$, respectively, the changes from silence to $o_2$ and $o_{12}$ cost
$(\kappa-F(t))u$. The other decisive option costs
$(G(t)-\kappa)c_3>0$ relative to silence. Contradictory advice costs more
than either decisive option. These observations prove optimality of
$B_\theta$.

A coherent minimizer must give $o_{23}$ almost everywhere on $P_1$, so its
credal set lies on $[e_2,e_3]$: any member with $m_1>0$ would prevent this
advice on an open subset of $P_1$. It must contain $e_2$, since otherwise
it gives $o_3$ on a positive part of $P_3$, where that option is never
optimal. Hence it is $B_s$. The two weak inequalities, together with
weighted optimality and $\kappa>0$, imply equality in both losses.
Finally $s\mapsto L_2(B_s)$ is strictly increasing, so $s=\theta$.
\end{proof}

\begin{theorem}[Dense thresholds]\label{thm:dense}
For every finite $\theta>0$ there are good bases $C_j=\Dscr_{K_j}$ with
$K_j\to K_\theta$ in Hausdorff distance. In particular, their thresholds
$\tau_j$ tend to $\theta$.
\end{theorem}
\begin{proof}
We first make a small improvement of $B_\theta$ which cannot be dominated
by any nearby model on the edge $[e_2,e_3]$.

Choose $a_j\uparrow\theta$ and $b_j\downarrow\theta$ with
\[
 u(P_2\cap\{\theta<t<b_j\})
 =u(P_3\cap\{a_j<t<\theta\})=z_j>0.
\]
This is possible because both weighted ratio distributions are atomless
and positive on every interval. Denote these sets by $A_j$ and $T_j$.
Choose a non-null $Q_j\subset P_1$ so small that
\begin{align}
 c_2(Q_j)&<c_2(P_3\cap\{\theta<t<b_j\})+c_2(T_j),\label{eq:q2}\\
 c_3(Q_j)&<u((P_2\cup P_3)\cap\{a_j<t<\theta\}),\label{eq:q3}
\end{align}
and $\mu(Q_j)\to0$. Let $X_j$ be $B_\theta$ with its debt silenced on
$Q_j\cup A_j$, and with credit added on $T_j$. Then
\begin{align*}
 L_2(X_j)-L_2(B_\theta)&=c_2(Q_j)+c_2(A_j)-c_2(T_j),\\
 L_3(X_j)-L_3(B_\theta)&=c_3(Q_j)-u(A_j)+u(T_j)=c_3(Q_j).
\end{align*}
Both differences tend to zero.

No $B_s$ can dominate $X_j$. If $a_j\le s\le b_j$, $X_j$ is a pointwise
world-1 improvement of $B_s$, strict on $Q_j$. If $s>b_j$, then
\[
 L_2(B_s)-L_2(B_\theta)
 \ge c_2(A_j)+c_2(P_3\cap\{\theta<t<b_j\})
 >L_2(X_j)-L_2(B_\theta)
\]
by \eqref{eq:q2}. If $s<a_j$, then
\[
 L_3(B_s)-L_3(B_\theta)
 \ge u((P_2\cup P_3)\cap\{a_j<t<\theta\})
 >L_3(X_j)-L_3(B_\theta)
\]
by \eqref{eq:q3}.

By KKM choose selected $C_j=\Dscr_{K_j}\preceq_iX_j$ at all worlds.
We claim $K_j\to K_\theta$. The nonempty compact convex subsets of
$\Delta$ are compact in Hausdorff distance. Along any convergent
subsequence, coherent advice converges almost everywhere, so both ordinary
losses converge. Its limit $C_*$ therefore satisfies
$L_k(C_*)\le L_k(B_\theta)$, $k=2,3$. Since $B_\theta$ is globally
optimal for $L_2+\kappa L_3$, $C_*$ is also a minimizer of that weighted
loss. Lemma~\ref{lem:bayes} forces
$C_*=B_\theta$, and hence $K_*=K_\theta$.
Here we use the standard facts that the zero set of a coherent lower
prevision is null (translate along the all-ones direction), and that
distinct compact convex credal sets give distinct advice on an open set
(strict separation followed by such a translation).

For large $j$, $\min_{K_j}m_2>0$, and $K_j$ is not contained in
$\{m_3=0\}$. If $K_j\subset\{m_1=0\}$, the weights lemma forces
$C_j=B_s$: an edge interval omitting both endpoints would recommend both
$o_2$ on $P_2$ and $o_3$ on $P_3$, and a credal set containing $e_3$ is
excluded by $\min m_2>0$. But all $B_s$ were excluded above.
By the face lemma, the selection therefore has full support. It gives $o_2$ on a positive part of
$P_2$, so the no-both lemma says it never gives $o_{13}$ there, modulo
null sets. This forces $e_2\in K_j$: otherwise $(1,-\eta,1)$ is strictly
desirable for sufficiently small $\eta>0$, and remains so on an open set.
Consequently $C_j$ is a good base. The continuous function $m_3/m_2$ on
a neighborhood of $K_\theta$ gives $\tau_j\to\theta$.
\end{proof}

\section{What dense thresholds already rule out}

\begin{corollary}\label{cor:flat}
Neither $F$ nor $G$ can be constant almost everywhere on a nonempty open
interval.
\end{corollary}
\begin{proof}
Suppose $F=\kappa$ on an interval. Theorem~\ref{thm:dense} supplies a good
base whose threshold lies strictly inside it. Choose a non-null debt set
$Y\subset P_2$ and a non-null silent set $Y'\subset P_3$, both close enough
to that threshold that $c_2=\kappa u$ on both. These are precisely the two
sets needed in IC's flat-threshold proof. That proof only uses constancy
on its chosen debt and silence sets; it does not require constancy on the
whole interval down to zero. Its four accepted changes and final guidance
worsening therefore give the forbidden thin-band comparisons. Exchange
worlds $2,3$ for $G$; the transformed function is $s\mapsto1/G(1/s)$.
\end{proof}

\begin{corollary}[Jumps everywhere]\label{cor:jumps}
Under the standing hypotheses, the discontinuities of $F$ are dense in
$(0,\infty)$, and so are those of $G$.
\end{corollary}
\begin{proof}
If $1/\rho$ is essentially unbounded on $P_1\cap\{y<\ell\}$ for every
$\ell>0$, IC's unbounded-ratio theorem already gives a contradiction.
Otherwise, for some $\ell,r>0$, $\rho\ge r$ on $P_1\cap\{y<\ell\}$.
At any good base with $F$ continuous at its threshold, IC's continuity
proof applies: its initial debt set can be chosen with
$y<\min(1,\ell)$, and the same uniform lower bound $r$ controls its cost.
Thus every good-base threshold must be a discontinuity of $F$.
Theorem~\ref{thm:dense} makes these thresholds dense. The mirror argument
gives the conclusion for $G$.
\end{proof}

This rules out much more than isolated jumps accumulating at zero or
infinity. It does not rule out a strictly increasing monotone function
whose countably many jumps are dense.

\section{Repeated debt reduction}

A jump need not prevent a chain from closing. It is enough that each round
reduce the remaining debt by a fixed proportion. The next lemma makes the
required inequality explicit.

\begin{lemma}[Local debt cycle]\label{lem:contract}
Let $C$ be a good base. Suppose $Y$ is a non-null debt set of $C$ in
$P_2\cap\{x\ge\ell\}$, and $Y_0'$ a non-null silent set in
$P_3\cap\{t<\tau\}$, with $0<\ell<\tau_2$. Suppose matching their
$u$-masses gives $c_2(Y')\ge\lambda c_2(Y_*)$, for some $0<\lambda\le1$.
Suppose there are non-null sets $M_0,R\subset P_1$, with $M_0$ silent
at $C$, and finite $\theta_0<b$, such that
\[
 M_0\subset\{y<\theta_0\},\qquad
 R\subset\{x<\ell,\ \theta_0<y<b\},\qquad
 \rho\le B\text{ on }M_0,\quad \rho\ge A>0\text{ on }R.
\]
If $(B/A)(1-\lambda)<1$, then $V$ is dominated.
\end{lemma}
\begin{proof}
Fix a sufficiently small non-null $S\subset P_3\cap\{t>\tau,y\ge b\}$.
Choose $\varepsilon>0$ so small that
\[
 q=\frac BA(1+\varepsilon)
       \big[(1-\lambda)(1+\varepsilon)+\varepsilon\big]<1.
\]
Start with a small debt set $Q_0\subset R$ and shrink $S$ further so that
$u(S)<c_3(Q_0)$. Competitiveness of $V$ gives
$V\preceq_1V+o_{12}(S)+o_{23}(Q_0)$: both other losses decrease.

Suppose the current debt is $Q_k\subset R$, of $c_2$-mass $x_k$.
The following round preserves $S$ and returns all auxiliary sets to silence.
\begin{enumerate}[label=\arabic*.,leftmargin=*,itemsep=3pt]
\item At base $E_2(\ell)$, move the debt from $Q_k$ to $Y_k\subset Y$ with
$c_2(Y_k)=(1+\varepsilon)x_k$. The change in $L_2$ is $-\varepsilon x_k$.
\item At the good base, silence $Y_k$, add credit on $Y_k'\subset Y_0'$
with $u(Y_k')=u(Y_k)$, and add debt on $M_k\subset M_0$ with
\[
 c_2(M_k)=m_k=(c_2(Y_k)-c_2(Y_k'))^++\varepsilon x_k.
\]
The changes satisfy $\Delta L_2\le-\varepsilon x_k$ and
$\Delta L_3=-c_3(M_k)<0$. Moreover
$m_k\le[(1-\lambda)(1+\varepsilon)+\varepsilon]x_k$.
\item Erase the added credit on $Y_k'$; this is a guidance worsening at
world $1$ and so preserves the accepted chain.
\item At base $E_3(\theta_0)$, move the debt from $M_k$ to
$Q_{k+1}\subset R$ with $c_3(Q_{k+1})=(1+\varepsilon)c_3(M_k)$.
This decreases $L_3$ strictly, and
\[
 x_{k+1}=c_2(Q_{k+1})
 \le\frac1A c_3(Q_{k+1})
 \le\frac BA(1+\varepsilon)m_k\le qx_k.
\]
\end{enumerate}
All sets exist by atomlessness if $Q_0$ is small enough. The same positive
reservoirs can be reused because they are silent again after each round.
The regions of the simultaneous changes are disjoint; in particular,
$S$ lies above $\tau$ in $P_3$ and $Y_k'$ lies below it.

For any sufficiently large $n$, take enough rounds that
$x_k<c_2(Z_n)/2$, where
$Z_n=P_1\cap\{b-1/n\le y<b,\ x\ge\ell\}$.
One $E_2(\ell)$ change moves the remaining debt into $D_n\subset Z_n$
with $c_2(D_n)>x_k$. Add debt on the rest of $B_n$ by (G).
This gives the thin-band comparison forbidden by the closing lemma,
for every sufficiently large $n$.
\end{proof}

\begin{lemma}[Flexible temporary credit]\label{lem:flexible}
Suppose $\rho\ge A>0$ on $P_1\cap\{x<h\}$. If a good base has threshold
$\tau$ satisfying
\[
 \frac1{F(\tau-)}-\frac1{F(\tau+)}<A,
 \tag{*}\label{eq:contraction}
\]
then $V$ is dominated.
\end{lemma}
\begin{proof}
Choose $0<\beta\le\alpha<\infty$ and non-null sets $Y$ in the base's
$P_2$ debt and $Y_0'$ in its $P_3$ silence, close to $\tau$ on opposite
sides, so that
\[
 F\le\alpha\text{ on }Y,\qquad F\ge\beta\text{ on }Y_0',\qquad
 1/\beta-1/\alpha<A.
\]
This is possible by the strict inequality in the hypothesis. Shrink $Y$
so that $x\ge\ell$ on it, with $0<\ell<\min(h,\tau_2)$.
Choose any bounded non-null silent reservoir $M_0\subset P_1$ of the base,
and then choose $\theta_0$ so large that $y<\theta_0$ on $M_0$.
There is no prescribed upper bound on $\theta_0$. Put
\[
 R=P_1\cap\{x<\ell,\ \theta_0<y<\theta_0+1\},\qquad b=\theta_0+2.
\]
Thus $\rho\ge A$ on $R$. No bound on $\rho$ is imposed on $M_0$.

Choose $\varepsilon>0$ so small that $\varepsilon<1/\alpha$ and, with
\[
 \gamma=1-\beta/\alpha+2\beta\varepsilon>0,\qquad
 q=\frac{(1+\varepsilon)^2\gamma}{A\beta},
\]
we have $q<1$. Start as in the local debt cycle with a sufficiently small
$Q_0\subset R$ and a fixed non-null
$S\subset P_3\cap\{t>\tau,y\ge b\}$, taking $u(S)<c_3(Q_0)$.

Suppose the current debt is $Q_k\subset R$, with $c_2$-mass $x_k$.
An $E_2(\ell)$ step moves it to $Y_k\subset Y$ with
$z_k=c_2(Y_k)=(1+\varepsilon)x_k$. Choose $M_k\subset M_0$ using the
atomless measure with density $c_2+\beta c_3$, so that
\[
 m_k+\beta r_k=\gamma z_k,\qquad
 m_k=c_2(M_k),\quad r_k=c_3(M_k)>0.
\]
This time choose the temporary credit $Y_k'\subset Y_0'$ with
\[
 u(Y_k')=u(Y_k)+r_k-\varepsilon z_k.
\]
Its prescribed mass is positive because $u(Y_k)\ge z_k/\alpha$.
At the good base, silence $Y_k$, add credit on $Y_k'$, and add debt on
$M_k$. Then
\begin{align*}
 \Delta L_3&=-u(Y_k)+u(Y_k')-r_k=-\varepsilon z_k<0,\\
 \Delta L_2&=z_k-c_2(Y_k')-m_k\\
 &\le z_k-\beta\bigl(u(Y_k)+r_k-\varepsilon z_k\bigr)-m_k\\
 &\le (1-\beta/\alpha)z_k-(m_k+\beta r_k)+\beta\varepsilon z_k
   =-\beta\varepsilon z_k<0.
\end{align*}
Erase the temporary credit by (G). An $E_3(\theta_0)$ step moves the debt
from $M_k$ to $Q_{k+1}\subset R$ with
$c_3(Q_{k+1})=(1+\varepsilon)r_k$. Consequently
\[
 c_2(Q_{k+1})\le\frac{(1+\varepsilon)r_k}{A}
 \le\frac{(1+\varepsilon)\gamma z_k}{A\beta}=q x_k.
\]
All prescribed masses are bounded by constants times $x_k$: on $Y$,
$F$ has a positive lower bound, and $r_k\le\gamma z_k/\beta$.
Thus all reservoirs suffice for a small enough initial $Q_0$, and are
silent again after each round. The same finite-band closing argument as
in Lemma~\ref{lem:contract} gives the contradiction.
\end{proof}

\begin{theorem}[A bound on one strip suffices]\label{thm:bounded}
Suppose worlds $2,3$ have ordinary additive rankings as above, and world $1$
is any preorder satisfying (G) and (IB). If (S) holds and, for some
$h>0$, either
\[
 \frac{c_3}{c_2}\ge A>0\quad\text{on }P_1\cap\{t_{12}<h\},
 \qquad\text{or}\qquad
 \frac{c_3}{c_2}\le B<\infty\quad\text{on }P_1\cap\{t_{13}<h\},
\]
then $V$ is dominated. This remains true when the jumps of $F$ and $G$
are dense.
\end{theorem}
\begin{proof}
Assume $V$ undominated. Choose any continuity point $\theta>0$ of $F$;
one exists because a monotone function has at most countably many
discontinuities. By Theorem~\ref{thm:dense} there are good bases with
$\tau_j\to\theta$. Monotonicity and continuity at $\theta$ give
\[
 F(\tau_j-)\longrightarrow F(\theta)>0,
 \qquad F(\tau_j+)\longrightarrow F(\theta).
\]
Thus \eqref{eq:contraction} holds for large $j$, contradicting
Lemma~\ref{lem:flexible}. The other strip follows by exchanging
worlds $2,3$, which replaces $\rho$ by its reciprocal.
\end{proof}

\begin{corollary}\label{cor:onesided}
A global finite upper bound on $\rho$, or a global positive lower bound
on $\rho$, suffices for domination of $V$.
\end{corollary}
\begin{proof}
Either global bound implies the corresponding strip bound in
Theorem~\ref{thm:bounded}.
\end{proof}

\section{Separable ratios}

The local cycle of Lemma~\ref{lem:contract} only needs an upper bound
for $\rho$ on its auxiliary silent set $M_0$, and a lower bound on its
returning debt reservoir $R$. Choosing these reservoirs explicitly handles
some unbounded ratios.

\begin{lemma}\label{lem:factor}
Suppose $\rho=f(x)/h(y)$ on $P_1$, where $f,h$ are finite positive
nondecreasing functions. At any good base put $a=\tau_2$ and
$\lambda=F(\tau/2)/\sup F>0$. If
\[
 \frac{f(a+)}{f(a-)}(1-\lambda)<1,
 \tag{**}\label{eq:factor}
\]
then $V$ is dominated. In particular this holds if $f$ is continuous at $a$.
\end{lemma}
\begin{proof}
Choose $x_-<\ell<a<x_+$ sufficiently close to $a$ that
$(f(x_+)/f(x_-))(1-\lambda)<1$. The base gives debt on a non-null part
of $P_2\cap\{x\ge\ell\}$: take $t$ sufficiently large in the shape
inequality. Choose $Y$ there, and choose a non-null silent
$Y_0'\subset P_3\cap\{\tau/2<t<\tau\}$. Matching their $u$-masses gives
$c_2(Y')\ge\lambda c_2(Y_*)$ as in the debt cycle.

There is a credal member attaining $m_2/m_1=a$, with $m_1>0$.
For $x$ close to $x_+>a$ and all sufficiently large $y$, its scalar
product with $(g_1,-|g_2|,-|g_3|)$ is positive. Since $e_2$ gives a
negative scalar product, this is an open silent region of the base.
Choose a sufficiently large continuity point $\theta_0$ of $h$.
Take a small silent reservoir $M_0$ with $x<x_+$ and
$\theta_0-\delta<y<\theta_0$, and a returning reservoir $R$ with
$x_-<x<\ell$ and $\theta_0<y<\theta_0+\delta$. Then
\[
 \rho\le B=\frac{f(x_+)}{h(\theta_0-\delta)}\quad\hbox{on }M_0,
 \qquad
 \rho\ge A=\frac{f(x_-)}{h(\theta_0+\delta)}\quad\hbox{on }R.
\]
Continuity of $h$ allows $\delta$ so small that $(B/A)(1-\lambda)<1$.
Choose $\varepsilon>0$ as in Lemma~\ref{lem:contract}. Its same four
changes now contract the debt by a fixed factor less than one, and its
closing argument gives the contradiction. No continuity of $F$ at $\tau$
has been used.
\end{proof}

\begin{theorem}\label{thm:factor}
For a separable ratio as in Lemma~\ref{lem:factor}, $V$ is dominated if
there are $J<\infty$ and $a_0<\infty$ such that
\[
 \frac{f(a+)}{f(a-)}\le J\quad\text{for all }a>a_0,
 \qquad\text{or}\qquad
 \frac{h(a+)}{h(a-)}\le J\quad\text{for all }a>a_0.
\]
In particular, continuity of either factor suffices, however fast it grows.
\end{theorem}
\begin{proof}
Suppose the bound holds for $f$. Write $d=\sup F$ and choose $\theta$ so
large that $J(1-F(\theta/4)/d)<1$; this is possible because $F(t)\to d$.
Take the good bases of Theorem~\ref{thm:dense} at this $\theta$.
Their thresholds $\tau_j\to\theta$, while $a_j=\tau_{2,j}\to\infty$:
their credal sets approach $K_\theta\subset\{m_1=0\}$, with $m_2$
uniformly bounded away from zero. For large $j$, $a_j>a_0$ and
$F(\tau_j/2)/d\ge F(\theta/4)/d$, so \eqref{eq:factor} holds.
Lemma~\ref{lem:factor} gives the contradiction. Exchange worlds $2,3$
for the assertion about $h$.
\end{proof}

This includes $\rho=(1+x)/(1+y)$, which escapes all the strip bounds above.
The theorem assumes separability explicitly. The necessity of separability
has now been written out in \emph{Edge Cycles and Separability}; its proof
does not assume completeness.

\section{The remaining gap}

The general conjecture remains open. CCF settles the case of two finitely
additive one-sided extensions when the remaining world is complete. The
unrestricted case of two ordinary additive worlds remains open only when
the remaining preorder is allowed to be incomplete. The new conclusions are necessary conditions,
not a construction of a counterexample.

In the latter case a counterexample now needs dense jumps of both $F$ and
$G$. Moreover, $\rho$ must be unbounded above on every strip $y<h$,
and $1/\rho$ must be unbounded above on every strip $x<h$. These are
additional to IC's bounds in the opposite directions. Indeed IC forces a
positive lower bound on $\rho$ in some strip $y<h$; a finite upper bound
on any such strip would invoke Theorem~\ref{thm:bounded}. The mirror
argument applies to $x<h$. These requirements are consistent at the level
of ratio functions. For example,
\[
 \rho(x,y)=\frac{1+x}{1+y}
\]
is bounded above on every strip $x<\ell$, and its reciprocal is bounded
above on every strip $y<\ell$, but $\rho$ is unbounded on every strip $y<h$ and $1/\rho$ is unbounded
on every strip $x<h$. This particular ratio is nevertheless excluded by
Theorem~\ref{thm:factor}. A separable counterexample must instead have
multiplicative jump ratios of both factors that are unbounded at
arbitrarily large arguments. More precisely, at every good-base parameter $a=\tau_2$ it must
fail \eqref{eq:factor}. These are necessary conditions only, not an
existence argument for a ranking.

In Theorem~\ref{thm:dense}, $K_j\to K_\theta$ also forces
$\tau_{2,j}\to\infty$. Consequently the auxiliary silent debt sets used
in the contraction proof can move to regions where $\rho$ is arbitrarily
large, while the returning debt sets can have arbitrarily small $\rho$.
Without a strip bound or a bound on the factors' multiplicative jumps, the
exchange cost may grow faster than the jump of $F$ shrinks. The proofs
above do not control that competition of rates.
Nor do these arguments establish path-closedness, or an additive extension, for three
arbitrary ordinal rankings.
\end{document}
